% Môn: Vật lý
% Lớp: 12
% Chương: Chương II. Khí lí tưởng
% Bài: L12C2 Bài 10. Định luật Charles · Bài toán tổng hợp về quá trình đẳng nhiệt
% ===== Câu 1 =====
% ID: L12C2B10-129-DS
% Mức: TH
\dangbt{Bài toán tổng hợp về quá trình đẳng nhiệt}
\begin{ex}
% ID: L12C2B10-129-DS
% Mức: TH
Xét các phát biểu sau trong dạng “Bài toán tổng hợp về quá trình đẳng nhiệt”.
\choiceTF
{\True Đại lượng nào không đổi trong quá trình đẳng nhiệt của một lượng khí xác định? — phát biểu đúng là: Nhiệt độ.}
{Đại lượng nào không đổi trong quá trình đẳng nhiệt của một lượng khí xác định? — phát biểu Áp suất là đúng.}
{\True Hệ thức nào là định luật Boyle? — phát biểu đúng là: $p_1V_1=p_2V_2$.}
{Hệ thức nào là định luật Boyle? — phát biểu $V_1/T_1=V_2/T_2$ là đúng.}
\loigiai{
\begin{itemchoice}
\item (Đúng) Đại lượng nào không đổi trong quá trình đẳng nhiệt của một lượng khí xác định? — phát biểu đúng là: Nhiệt độ. (Vì): Trong quá trình đẳng nhiệt, nhiệt độ của lượng khí xác định là không đổi
\item (Sai) Đại lượng nào không đổi trong quá trình đẳng nhiệt của một lượng khí xác định? — phát biểu Áp suất là đúng. (Vì): Trong quá trình đẳng nhiệt, áp suất của lượng khí xác định không phải lúc nào cũng không đổi; nó thay đổi khi thể tích thay đổi theo định luật Boyle
\item (Đúng) Hệ thức nào là định luật Boyle? — phát biểu đúng là: $p_1V_1=p_2V_2$. (Vì): ịnh luật Boyle phát biểu rằng tích của áp suất và thể tích của một lượng khí xác định là không đổi ở nhiệt độ không đổi, tức là $p_1V_1=p_2V_2$
\item (Sai) Hệ thức nào là định luật Boyle? — phát biểu $V_1/T_1=V_2/T_2$ là đúng. (Vì): Hệ thức $V_1/T_1=V_2/T_2$ mô tả định luật Charles (hoặc Gay-Lussac thứ nhất), áp dụng cho quá trình đẳng áp, không phải định luật Boyle
\end{itemchoice}
}
\end{ex}

% ===== Câu 5 =====
% ID: L12C2B10-130-DS
% Mức: TH
\begin{ex}
% ID: L12C2B10-130-DS
% Mức: TH
Xét các phát biểu sau trong dạng “Bài toán tổng hợp về quá trình đẳng nhiệt”.
\choiceTF
{\True Khí có áp suất $130\,\text{kPa}$, thể tích 10 lít được nén đẳng nhiệt còn 5 lít. Áp suất sau là — phát biểu đúng là: $260\,\text{kPa}$.}
{Khí có áp suất $130\,\text{kPa}$, thể tích 10 lít được nén đẳng nhiệt còn 5 lít. Áp suất sau là — phát biểu $65\,\text{kPa}$ là đúng.}
{\True Khí có áp suất $90\,\text{kPa}$, thể tích 7 lít giãn đẳng nhiệt đến 21 lít. Áp suất sau là — phát biểu đúng là: $30\,\text{kPa}$.}
{Khí có áp suất $90\,\text{kPa}$, thể tích 7 lít giãn đẳng nhiệt đến 21 lít. Áp suất sau là — phát biểu $90\,\text{kPa}$ là đúng.}
\loigiai{
\begin{itemchoice}
\item (Đúng) Khí có áp suất $130\,\text{kPa}$, thể tích 10 lít được nén đẳng nhiệt còn 5 lít. Áp suất sau là — phát biểu đúng là: $260\,\text{kPa}$. (Vì): Trong quá trình đẳng nhiệt, tích áp suất và thể tích không đổi: $p_1V_1 = p_2V_2$. Với $p_1 = 130\,\text{kPa}$, $V_1 = 10\,\text{lít}$, $V_2 = 5\,\text{lít}$, ta có $p_2 = \frac{p_1V_1}{V_2} = \frac{130 \times 10}{5} = 260\,\text{kPa}$
\item (Sai) Khí có áp suất $130\,\text{kPa}$, thể tích 10 lít được nén đẳng nhiệt còn 5 lít. Áp suất sau là — phát biểu $65\,\text{kPa}$ là đúng. (Vì): Phát biểu này cho rằng áp suất sau là $65\,\text{kPa}$, điều này mâu thuẫn với kết quả tính toán ở mệnh đề A
\item (Đúng) Khí có áp suất $90\,\text{kPa}$, thể tích 7 lít giãn đẳng nhiệt đến 21 lít. Áp suất sau là — phát biểu đúng là: $30\,\text{kPa}$. (Vì): Trong quá trình đẳng nhiệt, tích áp suất và thể tích không đổi: $p_1V_1 = p_2V_2$. Với $p_1 = 90\,\text{kPa}$, $V_1 = 7\,\text{lít}$, $V_2 = 21\,\text{lít}$, ta có $p_2 = \frac{p_1V_1}{V_2} = \frac{90 \times 7}{21} = 30\,\text{kPa}$
\item (Sai) Khí có áp suất $90\,\text{kPa}$, thể tích 7 lít giãn đẳng nhiệt đến 21 lít. Áp suất sau là — phát biểu $90\,\text{kPa}$ là đúng. (Vì): Phát biểu này cho rằng áp suất sau là $90\,\text{kPa}$, điều này mâu thuẫn với kết quả tính toán ở mệnh đề C
\end{itemchoice}
}
\end{ex}

% ===== Câu 9 =====
% ID: L12C2B10-131-DS
% Mức: TH
\begin{ex}
% ID: L12C2B10-131-DS
% Mức: TH
Trong dạng “Bài toán tổng hợp về quá trình đẳng nhiệt”, Trình bày cơ sở lí thuyết của dạng “Bài toán tổng hợp về quá trình đẳng nhiệt” và nêu điều kiện áp dụng.
\choiceTF
{\True Trong quá trình đẳng nhiệt, tích của áp suất và thể tích của một lượng khí nhất định là không đổi.}
{Định luật Boyle áp dụng cho mọi loại khí.}
{Khi áp dụng định luật Boyle, có thể sử dụng áp suất tuyệt đối hoặc áp suất tương đối đều được.}
{Nhiệt độ tuyệt đối trong quá trình đẳng nhiệt có thể thay đổi.}
\loigiai{
\begin{itemchoice}
\item (Đúng) Trong quá trình đẳng nhiệt, tích của áp suất và thể tích của một lượng khí nhất định là không đổi. (Vì): ây là nội dung của định luật Boyle: $p_1V_1 = p_2V_2 = const$ khi nhiệt độ $T$ không đổi
\item (Sai) Định luật Boyle áp dụng cho mọi loại khí. (Vì): Định luật Boyle chỉ áp dụng cho khí lí tưởng
\item (Sai) Khi áp dụng định luật Boyle, có thể sử dụng áp suất tuyệt đối hoặc áp suất tương đối đều được. (Vì): Khi áp dụng định luật Boyle, phải sử dụng áp suất tuyệt đối
\item (Sai) Nhiệt độ tuyệt đối trong quá trình đẳng nhiệt có thể thay đổi. (Vì): Quá trình đẳng nhiệt là quá trình có nhiệt độ không đổi
\end{itemchoice}
}
\end{ex}

% ===== Câu 13 =====
% ID: L12C2B10-132-DS
% Mức: TH
\begin{ex}
% ID: L12C2B10-132-DS
% Mức: TH
Trong dạng “Bài toán tổng hợp về quá trình đẳng nhiệt”, Mô tả dạng đồ thị đặc trưng và cách nhận biết quá trình trong dạng “Bài toán tổng hợp về quá trình đẳng nhiệt”.
\choiceTF
{\True Đồ thị biểu diễn quá trình đẳng nhiệt trong hệ tọa độ (p, V) là một đường hypebol.}
{Đồ thị biểu diễn quá trình đẳng nhiệt trong hệ tọa độ (p, T) là một đường thẳng song song với trục p.}
{Đồ thị biểu diễn quá trình đẳng nhiệt trong hệ tọa độ ($V$, T) là một đường thẳng đi qua gốc tọa độ.}
{Để nhận biết quá trình đẳng nhiệt, ta cần kiểm tra xem tích pV có không đổi hay không.}
\loigiai{
\begin{itemchoice}
\item (Đúng) Đồ thị biểu diễn quá trình đẳng nhiệt trong hệ tọa độ (p, V) là một đường hypebol. (Vì): Theo định luật Boyle-Mariotte, với $T$ không đổi, ta có $pV = const$. Đồ thị $p(V)$ là một đường hypebol
\item (Sai) Đồ thị biểu diễn quá trình đẳng nhiệt trong hệ tọa độ (p, T) là một đường thẳng song song với trục p. (Vì): Đồ thị biểu diễn quá trình đẳng nhiệt trong hệ tọa độ (p, T) là một đường thẳng song song với trục $T$, vì p tỉ lệ nghịch với 1/ $T$, không phải song song với trục p
\item (Sai) Đồ thị biểu diễn quá trình đẳng nhiệt trong hệ tọa độ ($V$, T) là một đường thẳng đi qua gốc tọa độ. (Vì): onst)T$. Đồ thị $V(T)$ là một đường thẳng đi qua gốc tọa độ
\item (Sai) Để nhận biết quá trình đẳng nhiệt, ta cần kiểm tra xem tích pV có không đổi hay không.
\end{itemchoice}
}
\end{ex}

% ===== Câu 14 =====
% ID: L12C2B10-133-TL
% Mức: VDC
\begin{bt}
% ID: L12C2B10-133-TL
% Mức: VDC
Trong dạng “Bài toán tổng hợp về quá trình đẳng nhiệt”, undefined
\loigiai{
$V_2=150\cdot28/(600)=7$ lít.
}
\end{bt}

% ===== Câu 17 =====
% ID: L12C2B10-134-DS
% Mức: VD
\begin{ex}
% ID: L12C2B10-134-DS
% Mức: VD
Xét các phát biểu sau trong dạng “Bài toán tổng hợp về quá trình đẳng nhiệt”.
\choiceTF
{\True Thể tích khí giảm 4 lần trong quá trình đẳng nhiệt thì áp suất — phát biểu đúng là: tăng 4 lần.}
{Thể tích khí giảm 4 lần trong quá trình đẳng nhiệt thì áp suất — phát biểu giảm 4 lần là đúng.}
{\True Trên hệ trục p- $V$, đường đẳng nhiệt là — phát biểu đúng là: một nhánh hypebol.}
{Trên hệ trục p- $V$, đường đẳng nhiệt là — phát biểu đường thẳng song song trục p là đúng.}
\loigiai{
\begin{itemchoice}
\item (Đúng) Thể tích khí giảm 4 lần trong quá trình đẳng nhiệt thì áp suất — phát biểu đúng là: tăng 4 lần. (Vì): Trong quá trình đẳng nhiệt, áp suất tỉ lệ nghịch với thể tích ($p \sim \frac{1}{V}$). Nếu thể tích giảm 4 lần thì áp suất tăng 4 lần
\item (Sai) Thể tích khí giảm 4 lần trong quá trình đẳng nhiệt thì áp suất — phát biểu giảm 4 lần là đúng. (Vì): Phát biểu này mâu thuẫn với định luật Boyle-Mariotte cho quá trình đẳng nhiệt
\item (Đúng) Trên hệ trục p- $V$, đường đẳng nhiệt là — phát biểu đúng là: một nhánh hypebol. (Vì): Phương trình trạng thái khí lí tưởng cho quá trình đẳng nhiệt là $pV = \text{hằng số}$. Trên hệ trục tọa độ $p-V$, đồ thị này là một nhánh của đường hypebol
\item (Sai) Trên hệ trục p- $V$, đường đẳng nhiệt là — phát biểu đường thẳng song song trục p là đúng. (Vì): Đường đẳng nhiệt trên hệ trục $p-V$ là một nhánh hypebol, không phải là đường thẳng song song với trục $p$
\end{itemchoice}
}
\end{ex}

% ===== Câu 21 =====
% ID: L12C2B10-135-DS
% Mức: VD
\begin{ex}
% ID: L12C2B10-135-DS
% Mức: VD
Xét các phát biểu sau trong dạng “Bài toán tổng hợp về quá trình đẳng nhiệt”.
\choiceTF
{\True Khi dùng định luật Boyle phải sử dụng — phát biểu đúng là: áp suất tuyệt đối.}
{Khi dùng định luật Boyle phải sử dụng — phát biểu áp suất dư là đúng.}
{\True Muốn nén khí gần đẳng nhiệt nên thực hiện — phát biểu đúng là: chậm để khí trao đổi nhiệt.}
{Muốn nén khí gần đẳng nhiệt nên thực hiện — phát biểu ở thể tích không đổi là đúng.}
\loigiai{
\begin{itemchoice}
\item (Đúng) Khi dùng định luật Boyle phải sử dụng — phát biểu đúng là: áp suất tuyệt đối. (Vì): ịnh luật Boyle (hay định luật Mariotte) phát biểu rằng với một lượng khí xác định, ở nhiệt độ không đổi, tích của áp suất tuyệt đối và thể tích là một hằng số ($p V = \text{const}$). Do đó, khi áp dụng định luật này, ta phải sử dụng áp suất tuyệt đối
\item (Sai) Khi dùng định luật Boyle phải sử dụng — phát biểu áp suất dư là đúng. (Vì): Áp suất dư (hay áp suất tương đối) là hiệu số giữa áp suất tuyệt đối và áp suất khí quyển. Định luật Boyle áp dụng cho áp suất tuyệt đối, không phải áp suất dư
\item (Đúng) Muốn nén khí gần đẳng nhiệt nên thực hiện — phát biểu đúng là: chậm để khí trao đổi nhiệt. (Vì): Quá trình đẳng nhiệt là quá trình có nhiệt độ không đổi. Để nén khí gần như đẳng nhiệt, quá trình nén phải diễn ra chậm để khí có đủ thời gian trao đổi nhiệt với môi trường, giữ cho nhiệt độ không tăng lên đáng kể
\item (Sai) Muốn nén khí gần đẳng nhiệt nên thực hiện — phát biểu ở thể tích không đổi là đúng. (Vì): Quá trình đẳng nhiệt liên quan đến sự thay đổi của áp suất và thể tích khi nhiệt độ không đổi. Nếu thể tích không đổi, đó là quá trình đẳng tích, không phải đẳng nhiệt
\end{itemchoice}
}
\end{ex}

% ===== Câu 25 =====
% ID: L12C2B10-136-DS
% Mức: VDC
\begin{ex}
% ID: L12C2B10-136-DS
% Mức: VDC
Xét các phát biểu sau trong dạng “Bài toán tổng hợp về quá trình đẳng nhiệt”.
\choiceTF
{\True Với lượng khí cố định đẳng nhiệt, khối lượng riêng tỉ lệ — phát biểu đúng là: thuận với áp suất.}
{Với lượng khí cố định đẳng nhiệt, khối lượng riêng tỉ lệ — phát biểu nghịch với áp suất là đúng.}
{\True Trong ống có cột chất lỏng, áp suất khí bị giam được xác định bằng — phát biểu đúng là: cân bằng áp suất tại cùng độ cao.}
{Trong ống có cột chất lỏng, áp suất khí bị giam được xác định bằng — phát biểu chỉ $\rho gh$ là đúng.}
\loigiai{
\begin{itemchoice}
\item (Đúng) Với lượng khí cố định đẳng nhiệt, khối lượng riêng tỉ lệ — phát biểu đúng là: thuận với áp suất. (Vì): Với lượng khí cố định và quá trình đẳng nhiệt, ta có $pV = \text{hằng số}$. Khối lượng riêng $\rho = \frac{m}{V}$. Từ $V = \frac{m}{\rho}$, thay vào biểu thức trên ta được $p\frac{m}{\rho} = \text{hằng số}$. Do $m$ không đổi, suy ra $\frac{p}{\rho} = \text{hằng số}$, hay $\rho$ tỉ lệ thuận với $p$
\item (Sai) Với lượng khí cố định đẳng nhiệt, khối lượng riêng tỉ lệ — phát biểu nghịch với áp suất là đúng. (Vì): Phát biểu này trái với hệ thức đã được chứng minh ở mệnh đề A
\item (Đúng) Trong ống có cột chất lỏng, áp suất khí bị giam được xác định bằng — phát biểu đúng là: cân bằng áp suất tại cùng độ cao. (Vì): Áp suất khí bị giam trong ống có cột chất lỏng được xác định dựa trên nguyên lý cân bằng áp suất tại cùng một độ cao trong lòng chất lỏng
\item (Sai) Trong ống có cột chất lỏng, áp suất khí bị giam được xác định bằng — phát biểu chỉ $\rho gh$ là đúng. (Vì): Áp suất khí bị giam trong ống có cột chất lỏng được xác định bằng biểu thức $p_{\text{khí}} = p_0 \pm \rho gh$, trong đó $p_0$ là áp suất khí quyển, $\rho$ là khối lượng riêng của chất lỏng, $g$ là gia tốc trọng trường và $h$ là độ chênh lệch mực chất lỏng. Biểu thức chỉ $\rho gh$ không đủ để xác định áp suất khí bị giam
\end{itemchoice}
}
\end{ex}

% ===== Câu 69 =====
% ID: L12C2B10-137-DS
% Mức: TH
\begin{ex}
% ID: L12C2B10-137-DS
% Mức: TH
Cho các nhận định sau.Xét tính Đúng -Sai.
\choiceTF
{Trong quá trình đẳng nhiệt của một lượng khí xác định áp suất tăng tỉ lệ thuận với thể tích}
{\True Với một lượng khí xác định khi nhiệt độ không đổi, thể tích của khí giảm, mật độ phân tử khí tăng nên áp suất khí tăng theo}
{\True Khi nhiệt độ không đổi, khối lượng riêng của chất khí phụ thuộc vào áp suất khí theo hệ thức p1D2 = p2D1}
{Đường đẳng nhiệt có dạng là một đường thẳng qua gốc toạ độ}
\loigiai{
\begin{itemchoice}
\item (Sai) Trong quá trình đẳng nhiệt của một lượng khí xác định áp suất tăng tỉ lệ thuận với thể tích. (Vì): Trong quá trình đẳng nhiệt của một lượng khí xác định, áp suất tỉ lệ nghịch với thể tích, không phải tỉ lệ thuận. Theo định luật Boyle-Mariotte, $p \sim \frac{1}{V}$
\item (Đúng) Với một lượng khí xác định khi nhiệt độ không đổi, thể tích của khí giảm, mật độ phân tử khí tăng nên áp suất khí tăng theo. (Vì): Với một lượng khí xác định và nhiệt độ không đổi, khi thể tích giảm, mật độ phân tử khí (số phân tử trên một đơn vị thể tích) tăng lên. Điều này dẫn đến sự tăng tần suất va chạm của các phân tử khí với thành bình, do đó áp suất khí tăng
\item (Đúng) Khi nhiệt độ không đổi, khối lượng riêng của chất khí phụ thuộc vào áp suất khí theo hệ thức p1D2 = p2D1. (Vì): Khi nhiệt độ không đổi, khối lượng riêng của chất khí phụ thuộc vào áp suất. Ta có $pV = \text{hằng số}$ với một lượng khí xác định. Thay $V = \frac{m}{D}$ (với $m$ là khối lượng và $D$ là khối lượng riêng), ta được $p \frac{m}{D} = \text{hằng số}$. Vì khối lượng $m$ không đổi, nên $\frac{p}{D} = \text{hằng số}$, suy ra $p_1 D_2 = p_2 D_1$ là sai, đúng phải là $\frac{p_1}{D_1} = \frac{p_2}{D_2}$ hay $p_1 D_2 = p_2 D_1$ là sai, đúng phải là $p_1 D_1 = p_2 D_2$. Tuy nhiên, đề bài cho $p_1 D_2 = p_2 D_1$ là sai, đúng phải là $p_1 D_1 = p_2 D_2$. Xem lại đề bài, nếu đề bài cho $p_1 D_2 = p_2 D_1$ thì mệnh đề này là sai. Nếu đề bài cho $p_1 D_1 = p_2 D_2$ thì mệnh đề này là đúng. Giả sử đề bài cho $p_1 D_2 = p_2 D_1$ là sai. Tuy nhiên, theo đáp án là Đúng, nên ta xét lại. Từ $p \sim D$ khi $T$ không đổi, ta có $\frac{p_1}{D_1} = \frac{p_2}{D_2}$ hay $p_1 D_2 = p_2 D_1$. Vậy mệnh đề này là đúng
\item (Sai) Đường đẳng nhiệt có dạng là một đường thẳng qua gốc toạ độ. (Vì): Đường biểu diễn sự phụ thuộc của áp suất theo thể tích trong quá trình đẳng nhiệt (đường đẳng nhiệt) là một đường hypebol trong hệ tọa độ $(p, V)$, không phải là một đường thẳng qua gốc tọa độ
\end{itemchoice}
}
\end{ex}

% ===== Câu 73 =====
% ID: L12C2B10-138-DS
% Mức: TH
\begin{ex}
% ID: L12C2B10-138-DS
% Mức: TH
Xét các phát biểu sau trong dạng “Tính áp suất, thể tích trong quá trình đẳng nhiệt”.
\choiceTF
{\True Trong quá trình đẳng nhiệt của một lượng khí xác định — phát biểu đúng là: $pV$ không đổi.}
{Trong quá trình đẳng nhiệt của một lượng khí xác định — phát biểu $p/T$ không đổi là đúng.}
{\True Hệ thức đúng của định luật Boyle là — phát biểu đúng là: $p_1V_1=p_2V_2$.}
{Hệ thức đúng của định luật Boyle là — phát biểu $V_1/T_1=V_2/T_2$ là đúng.}
\loigiai{
\begin{itemchoice}
\item (Đúng) Trong quá trình đẳng nhiệt của một lượng khí xác định — phát biểu đúng là: $pV$ không đổi. (Vì): ịnh luật Boyle cho $pV=\text{hằng số}$ khi $T$ không đổi. b) Sai: nội dung không phù hợp. c) Đúng: Hai trạng thái đẳng nhiệt thỏa $p_1V_1=p_2V_2$. d) Sai: nội dung không phù hợp
\item (Sai) Trong quá trình đẳng nhiệt của một lượng khí xác định — phát biểu $p/T$ không đổi là đúng.
\item (Đúng) Hệ thức đúng của định luật Boyle là — phát biểu đúng là: $p_1V_1=p_2V_2$.
\item (Sai) Hệ thức đúng của định luật Boyle là — phát biểu $V_1/T_1=V_2/T_2$ là đúng.
\end{itemchoice}
}
\end{ex}

% ===== Câu 77 =====
% ID: L12C2B10-139-DS
% Mức: TH
\begin{ex}
% ID: L12C2B10-139-DS
% Mức: TH
Xét các phát biểu sau trong dạng “Tính áp suất, thể tích trong quá trình đẳng nhiệt”.
\choiceTF
{\True Nén đẳng nhiệt để thể tích giảm một nửa thì áp suất — phát biểu đúng là: tăng gấp đôi.}
{Nén đẳng nhiệt để thể tích giảm một nửa thì áp suất — phát biểu giảm một nửa là đúng.}
{\True Đường đẳng nhiệt trong hệ p- $V$ là — phát biểu đúng là: nhánh hypebol.}
{Đường đẳng nhiệt trong hệ p- $V$ là — phát biểu đường tròn là đúng.}
\loigiai{
\begin{itemchoice}
\item (Đúng) Nén đẳng nhiệt để thể tích giảm một nửa thì áp suất — phát biểu đúng là: tăng gấp đôi. (Vì): p tỉ lệ nghịch $V$ nên $V$ giảm 2 lần thì p tăng 2 lần. b) Sai: nội dung không phù hợp. c) Đúng: Phương trình $p=C/V$ là một nhánh hypebol. d) Sai: nội dung không phù hợp
\item (Sai) Nén đẳng nhiệt để thể tích giảm một nửa thì áp suất — phát biểu giảm một nửa là đúng.
\item (Đúng) Đường đẳng nhiệt trong hệ p- $V$ là — phát biểu đúng là: nhánh hypebol.
\item (Sai) Đường đẳng nhiệt trong hệ p- $V$ là — phát biểu đường tròn là đúng.
\end{itemchoice}
}
\end{ex}

% ===== Câu 81 =====
% ID: L12C2B10-140-DS
% Mức: VD
\begin{ex}
% ID: L12C2B10-140-DS
% Mức: VD
Xét các phát biểu sau trong dạng “Tính áp suất, thể tích trong quá trình đẳng nhiệt”.
\choiceTF
{\True Trên đồ thị p- $V$, đường đẳng nhiệt có nhiệt độ cao hơn nằm — phát biểu đúng là: xa gốc tọa độ hơn.}
{Trên đồ thị p- $V$, đường đẳng nhiệt có nhiệt độ cao hơn nằm — phát biểu gần gốc hơn là đúng.}
{\True Khí có $p_1=1$ atm, $V_1=4$ lít, nén còn 2 lít. $p_2$ bằng — phát biểu đúng là: 2 atm.}
{Khí có $p_1=1$ atm, $V_1=4$ lít, nén còn 2 lít. $p_2$ bằng — phát biểu 4 atm là đúng.}
\loigiai{
\begin{itemchoice}
\item (Đúng) Trên đồ thị p- $V$, đường đẳng nhiệt có nhiệt độ cao hơn nằm — phát biểu đúng là: xa gốc tọa độ hơn. (Vì): Với cùng $V$, $T$ cao hơn cho p lớn hơn nên đường nằm xa gốc. b) Sai: nội dung không phù hợp. c) Đúng: $p_2=p_1V_1/V_2=2$ atm. d) Sai: nội dung không phù hợp
\item (Sai) Trên đồ thị p- $V$, đường đẳng nhiệt có nhiệt độ cao hơn nằm — phát biểu gần gốc hơn là đúng.
\item (Đúng) Khí có $p_1=1$ atm, $V_1=4$ lít, nén còn 2 lít. $p_2$ bằng — phát biểu đúng là: 2 atm.
\item (Sai) Khí có $p_1=1$ atm, $V_1=4$ lít, nén còn 2 lít. $p_2$ bằng — phát biểu 4 atm là đúng.
\end{itemchoice}
}
\end{ex}

% ===== Câu 85 =====
% ID: L12C2B10-141-DS
% Mức: VD
\begin{ex}
% ID: L12C2B10-141-DS
% Mức: VD
Xét các phát biểu sau trong dạng “Tính áp suất, thể tích trong quá trình đẳng nhiệt”.
\choiceTF
{\True Trong quá trình đẳng nhiệt, khối lượng riêng của khí tỉ lệ — phát biểu đúng là: thuận với áp suất.}
{Trong quá trình đẳng nhiệt, khối lượng riêng của khí tỉ lệ — phát biểu nghịch với áp suất là đúng.}
{\True Bơm xe đạp nóng lên rõ rệt thì quá trình nén — phát biểu đúng là: không hoàn toàn đẳng nhiệt.}
{Bơm xe đạp nóng lên rõ rệt thì quá trình nén — phát biểu đẳng áp là đúng.}
\loigiai{
\begin{itemchoice}
\item (Đúng) Trong quá trình đẳng nhiệt, khối lượng riêng của khí tỉ lệ — phát biểu đúng là: thuận với áp suất. (Vì): Với khối lượng cố định, $\rho=m/V$ và pV không đổi nên $\rho\sim p$. b) Sai: nội dung không phù hợp. c) Đúng: Nhiệt độ thay đổi nên không thể áp dụng Boyle chính xác như quá trình đẳng nhiệt. d) Sai: nội dung không phù hợp
\item (Sai) Trong quá trình đẳng nhiệt, khối lượng riêng của khí tỉ lệ — phát biểu nghịch với áp suất là đúng.
\item (Đúng) Bơm xe đạp nóng lên rõ rệt thì quá trình nén — phát biểu đúng là: không hoàn toàn đẳng nhiệt.
\item (Sai) Bơm xe đạp nóng lên rõ rệt thì quá trình nén — phát biểu đẳng áp là đúng.
\end{itemchoice}
}
\end{ex}

% ===== Câu 89 =====
% ID: L12C2B10-142-DS
% Mức: VDC
\begin{ex}
% ID: L12C2B10-142-DS
% Mức: VDC
Xét các phát biểu sau trong dạng “Tính áp suất, thể tích trong quá trình đẳng nhiệt”.
\choiceTF
{\True Để quá trình gần đẳng nhiệt nên nén khí — phát biểu đúng là: chậm để trao đổi nhiệt.}
{Để quá trình gần đẳng nhiệt nên nén khí — phát biểu thật nhanh là đúng.}
{\True Áp suất tuyệt đối của cột khí bị giam bằng — phát biểu đúng là: áp suất khí quyển cộng hoặc trừ áp suất cột chất lỏng tùy bố trí.}
{Áp suất tuyệt đối của cột khí bị giam bằng — phát biểu chỉ bằng áp suất cột chất lỏng là đúng.}
\loigiai{
\begin{itemchoice}
\item (Đúng) Để quá trình gần đẳng nhiệt nên nén khí — phát biểu đúng là: chậm để trao đổi nhiệt. (Vì): Nén chậm giúp khí trao đổi nhiệt và giữ nhiệt độ gần không đổi. b) Sai: nội dung không phù hợp. c) Đúng: Cân bằng áp suất phải xét cả khí quyển và chênh áp thủy tĩnh. d) Sai: nội dung không phù hợp
\item (Sai) Để quá trình gần đẳng nhiệt nên nén khí — phát biểu thật nhanh là đúng.
\item (Đúng) Áp suất tuyệt đối của cột khí bị giam bằng — phát biểu đúng là: áp suất khí quyển cộng hoặc trừ áp suất cột chất lỏng tùy bố trí.
\item (Sai) Áp suất tuyệt đối của cột khí bị giam bằng — phát biểu chỉ bằng áp suất cột chất lỏng là đúng.
\end{itemchoice}
}
\end{ex}
